\section{Discussion and Conclusion}
\label{sec:discussion}

\subsection{Interpretation}

The current evidence supports a narrow but important conclusion: India's agri-derivatives
suspension did not raise spot volatility in the banned commodities. The inherited
$+8$--$10\%$ claim is not supported by the cleaned district panel or by the current
counterfactual design.

The evidence does not support a broad claim that the suspension robustly reduced aggregate
spot volatility. Once food-cereal donors are added, the aggregate DiD is $-9.8\%$ and not
significant. The strongest lower-volatility evidence is concentrated in chana. Wheat, by
contrast, is not a clean treated unit: the current Synthetic DiD estimate is approximately
zero, and the post-2021 wheat market is heavily confounded by MSP/procurement, export
restrictions, stock limits and open-market sales.

The correct paper-level framing is therefore donor sensitivity and commodity heterogeneity.
A referee should see three claims, in this order: the inherited positive-volatility claim is
refuted; the aggregate lower-volatility estimate is modest and insignificant under a better
food-staple donor pool; chana remains the strongest candidate for a genuine lower-volatility
association.

\subsection{Mechanisms and confounding}

Several mechanisms could generate lower measured volatility without implying that the
suspension itself improved welfare. The ban was imposed after a hot food-price episode, so
mean reversion remains a threat. MSP and procurement can mechanically suppress measured
volatility in wheat and paddy. Export bans, stock limits, tariff changes and open-market
sales form a bundled policy environment that is difficult to separate from the futures
suspension.

This is why the paper should avoid a causal slogan. The defensible language is
``associated with'' lower volatility in chana and ``did not raise'' aggregate volatility.
It should not say the suspension caused a broad reduction in volatility.

\subsection{Cost side}

The welfare question is two-armed. Even if spot volatility did not rise, the suspension
removed price-discovery and hedging functions. H3 provides suggestive evidence that banned
spot markets became less spatially integrated relative to controls. H8 shows that direct
registered hedger turnover in NCDEX agri was very small (about 1.4\%), so the immediate
cost story is less ``farmers lost a large direct hedge'' and more ``the market lost liquidity,
reference pricing and integration.''

The inherited post-ban basis-widening claim remains impossible as stated because banned
commodities have no post-ban exchange-traded futures. The cost-side discovery and hedging
work must therefore be reframed as pre-ban price discovery, pre-ban hedging effectiveness,
or domestic-spot versus international-futures insulation.

\subsection{Limitations}

The main limitations are binding. The food/oilseed donor pull is incomplete: barley, maize
and jowar are complete; bajra is partial; ragi, groundnut, sesamum and sunflower are not in
the district panel. The oilseed donors matter especially for mustard and soybean. Few-cluster
inference remains coarse. The policy treatment is bundled with other commodity interventions.
Several literature sources remain paywalled/manual and need page-level verification before
submission. The paper bibliography is currently a compile stub, not a submission-grade
reference file.

\subsection{Conclusion}

The project has moved from a simple volatility-increase story to a more defensible and more
interesting result. The inherited claim that the suspension raised spot volatility is refuted.
The stronger clean-core claim that volatility fell by about 20\% does not survive a better
food-staple counterfactual as an aggregate significant result. The current result is that
aggregate spot volatility is estimated lower but insignificantly, with the credible signal
concentrated in chana and wheat confounded.

The next publishable step is mechanical and clear: finish the full food/oilseed donor pull,
rebuild the panel, rerun C1 estimators, add Honest-DiD and staggered-DiD robustness, and then
rewrite the final paper around the completed donor-sensitivity result.
